\documentclass[12pt]{article}
\usepackage[a4paper, margin=25mm]{geometry}
\usepackage{amsmath}
\usepackage{amsfonts}
\usepackage{mismath}
\title{Matrix differential equations}
\begin{document}
\maketitle
\appendix

\section{Matrix ODE derivation}
Note the following equation for a matrix variable $M$ and constant matrix $B$ of appropriate shapes:
\begin{equation*}
	\nabla_M\left(\tr(M B M^T)\right) = M(B^T + B).
\end{equation*}
In particular, for $B=I$, we find the gradient of the square of the Frobenius norm:
\begin{equation*}
	\nabla_M\left(\tr(M M^T)\right) = \nabla_M\left(\sum_{i,j} m_{ij}^2\right) = \nabla_M\left(\Vert M\Vert_F^2\right) = 2M.
\end{equation*}

Now let $X(t) \in \mathbb{R}^{m \times n}$ and define
\begin{equation*}
	\Phi(X) = \frac{1}{4} \Vert M(X) \Vert^2_F, \quad M(X) = X^T X - I.
\end{equation*}
By the matrix gradient chain rule, with a slight abuse of notation for the reparameterization of $\Phi$,
\begin{equation*}
	\nabla_X\Phi(X) = \frac{1}{4} \nabla_X M^T(X) \nabla_{M(X)} \Phi\left(M(X)\right) = X (X^T X - I).
\end{equation*}
Setting up an equation of gradient descent on potential $\Phi$ gives
\begin{equation*}
    \frac{dX}{dt} = -\nabla\Phi(X) = X(I - X^T X).
\end{equation*}

Note that we may produce an SVD factorization of $X = U \Sigma V^T.$ Afterwards,
\begin{align*}
    \frac{d}{dt}\left[U \Sigma V^T\right]
        &= U \Sigma V^T (I - V \Sigma^T U^T U \Sigma V^T) \\
        &= U \left( \Sigma \left( I - \Sigma^2 \right) \right) V^T.
\end{align*}
From this, it is clear that orthogonal factors $U, V$ stay fixed during evolution of $X(t)$, so the equation may be simplified to consider the evolution of the singular values:
\begin{equation*}
    \frac{d\sigma_i}{dt} = \sigma_i(1 - \sigma_i^2).
\end{equation*}
This ODE attains stable equilibrium points at $\sigma_i = \pm1$ and has an unstable equilibrium point at $\sigma_i = 0$. Therefore, any positive $\sigma_i$ will approach $1$ as $t\rightarrow\infty$, and the effect of the matrix differential equation is to take the input matrix to its polar factor $U V^T$.

Now consider that $\Phi$ is smooth and positive-definite on the matrix space, and see:
\begin{equation*}
    \frac{d}{dt} \Phi\left(X(t)\right) = \left(\nabla\Phi(X)\right)^T \frac{dX}{dt} = -\Vert\nabla\Phi(X)\Vert_F^2 \le 0.
\end{equation*}
Therefore $\Phi$ is Lyapunov on the problem space, and all zero solutions of the matrix ODE are stable. Note this does not contradict the singular value analysis as nonzero rank-deficient matrices are not zero solutions.

For the stability analysis, the gradient of the matrix equation RHS $-\nabla^2 \Phi$ is a quartix with the following components, using Einstein summation convention with index notation derivatives modified for matrices, i.e. $(\nabla f)_{ij} := \partial f/\partial X_{ij} := f_{,ij}$:
\begin{align*}
    -{(\nabla^2 \Phi)}_{abcd} &= -\Phi_{,ab,cd} = -\frac{1}{4} (M_{ij} M_{ij})_{,ab,cd} = -\frac{1}{2} (M_{ij} M_{ij,ab})_{,cd} \\
        &= -\frac{1}{2} (M_{ij,ab} M_{ij,cd} + M_{ij} M_{ij,ab,cd}).
\end{align*}
Now evaluate matrix derivatives of $M$:
\begin{align*}
    M_{ij,ab} &= (X^T X - I)_{ij,ab} = (X_{ki} X_{kj} - \delta_{ij})_{,ab} \\
        &= X_{ki,ab} X_{kj} + X_{ki} X_{kj,ab} = \delta_{ka} \delta_{ib} X_{kj} + X_{ki} \delta_{ka} \delta_{jb} \\
        &= \delta_{ib} X_{aj} + X_{ai} \delta_{jb}; \\
    M_{ij,ab,cd} &= \delta_{ib} X_{aj,cd} + X_{ai,cd} \delta_{jb} = \delta_{ib} \delta_{ac} \delta_{jd} +  \delta_{ac} \delta_{id} \delta_{jb} \\
        &= \delta_{ac} (\delta_{ib} \delta_{jd} + \delta_{id} \delta_{jb}).
\end{align*}
On the way, we may verify that the RHS expression derived earlier is correct:
\begin{align*}
    -{(\nabla \Phi)}_{ab} &= -\Phi_{,ab} = -\frac{1}{2} M_{ij} M_{ij,ab} = -\frac{1}{2} (X_{ki} X_{kj} - \delta_{ij})(\delta_{ib} X_{aj} + X_{ai} \delta_{jb}) \\
        &= \frac{1}{2} (\delta_{ij} \delta_{ib} X_{aj} - X_{ki} X_{kj} \delta_{ib} X_{aj} + \delta_{ij} X_{ai} \delta_{jb} - X_{ki} X_{kj} X_{ai} \delta_{jb}) \\
        &= X_{ab} - X_{aj} X_{kj} X_{kb} = X_{aj} (\delta_{jb} - X_{kj} X_{kb}) \\
        &= (X(I - X^T X))_{ab}.
\end{align*}
Therefore our matrix-valued matrix function Jacobian tensor may be unambiguously expressed using index notation:
\begin{align*}
    -{(\nabla^2 \Phi)}_{abcd} &= -\frac{1}{2} (M_{ij,ab} M_{ij,cd} + M_{ij} M_{ij,ab,cd}) \\
        &= -\frac{1}{2} ((\delta_{ib} X_{aj} + X_{ai} \delta_{jb}) (\delta_{id} X_{cj} + X_{ci} \delta_{jd}) + (X_{ki} X_{kj} - \delta_{ij}) \delta_{ac} (\delta_{ib} \delta_{jd} + \delta_{id} \delta_{jb})) \\
        &= -(X_{aj} \delta_{bd} X_{cj} + X_{ad} X_{cb}) + (\delta_{bd} - X_{kb} X_{kd}) \delta_{ac} \\
        &= \delta_{ac} \delta_{bd} -\delta_{bd} X_{aj} X_{cj} - \delta_{ac} X_{kb} X_{kd} - X_{ad} X_{cb}.
\end{align*}
We may vectorize $ab$ and $cd$ index pairs in order to conduct the usual eigenvector analysis.
\end{document}
